% F081: precise mathematical teaching diagram.
\begin{tikzpicture}[x=1cm,y=1cm,>=stealth,font=\small]
\draw[rounded corners=3pt,gray!40] (0,0) rectangle (14.2,8.6);
\begin{scope}[yshift=1cm]
\node at (7.1,7.0) {例3：上限为一个，再用异号端点保证存在};
\begin{scope}[shift={(4.5,3.3)},x=7cm,y=1.0cm]
\draw[->,gray!65] (-0.08,0)--(0.36,0) node[right] {$x-1$};
\draw[->,gray!65] (0,-1.8)--(0,2.4) node[above] {$y$};
\draw[blue!75!black,line width=1.15pt,,domain=-0.06:0.28,samples=160] plot (\x,{(\x+1)*(\x+1)*(\x+1)*(\x+1)*(\x+1)+(\x+1)-3});
\fill[blue!75!black] (0,-1) circle (1.8pt);\node[above left] at (0,-1) {$(1,-1)$};
\end{scope}
\node[] at (10.5,5.5) {$f^{\prime}(x)=5x^4+1>0$};
\node[] at (10.5,4.4) {全实数范围最多一根};
\node[] at (10.5,3.2) {$f(1)=-1,\quad f(2)=31$};
\node[] at (10.5,2.1) {连续保证至少一根};
\end{scope}
\node[text width=13cm,align=center] at (7.1,.9) {左图放大交点附近，横轴为 $x-1$；唯一零点位于 $(1,2)$。};
\end{tikzpicture}
